\documentclass[10pt,a4paper]{amsart}
\usepackage[margin=19mm]{geometry}
\usepackage{amsmath,amssymb,mathtools}
\usepackage[scr=rsfs,cal=euler]{mathalfa}
\usepackage{xfrac}
\usepackage[unicode,hidelinks,bookmarks=false]{hyperref}
\newcommand{\ie}{i.e.\ }
\newcommand{\cf}{cf.\ }
\newcommand{\resp}{respectively\ }
\newcommand{\Subsection}{\subsection}
\providecommand{\nobreakdash}{\nobreak\hbox{-}}
\setlength{\emergencystretch}{2em}
\multlinegap0pt
\usepackage{amssymb}
\usepackage{mathtools}
\usepackage{enumerate}
\usepackage[shortlabels]{enumitem}
\usepackage{tikz}
\usepackage{booktabs}
\newtheorem{definition}{Definition}
\newtheorem{lemma}{Lemma}
\newtheorem{theorem}{Theorem}
\newcommand{\Hnum}{H}                      % hypernumber universe
\newcommand{\R}{\mathbb{R}}
\newcommand{\Rp}{\mathbb{R}_{>0}}
\newcommand{\Snon}{S^{\times}}             % nonzero signs {+,-,Λ}
\newcommand{\iemb}{i}                      % embedding R -> H
\newcommand{\op}{\boxplus}                 % hyperaddition on H
\title{Three-sign cancellation hypernumber systems and associator curvature}
\author{Jaehwan Kim}
\date{Source: arXiv:2602.21207v1 (CC BY 4.0)}
\begin{document}
\maketitle

\noindent\emph{Source and licence.} Three-sign cancellation hypernumber systems and associator curvature. Version arXiv:2602.21207v1 (CC BY 4.0), distributed by arXiv under the Creative Commons Attribution 4.0 International licence, http://creativecommons.org/licenses/by/4.0/, as recorded in the arXiv licence metadata for this exact version and shown on the abstract page https://arxiv.org/abs/2602.21207v1. Full licence notice: \texttt{licenses/arxiv-2602.21207-CC-BY-4.0.txt}.\par\smallskip

\noindent\textit{Editorial scope.} Counted unit: the article's theorem thm:scalar-Rge0 (source0.tex lines 1875--1884). The article's complete original proof of it is delivered with it (source0.tex lines 1886--2062, one \texttt{proof} environment of 177 lines). No mathematics is written, repaired or re-derived: the statement and the proof are the article's own lines, in the article's own notation, and no retained line is substituted or re-flowed.  Retained uncounted as the source-local context the proof needs: lines 1413--1475; lines 1734--1760.  Nothing was omitted from any retained span, so the package carries no omission record.  No retained line contains a citation, so the wrapper carries no bibliography and no bibliography entry.  No retained line contains a figure environment, an image-inclusion command or drawing code, so no image is delivered and the package declares no asset dependency.  Editorial formatting only, in addition: an AMS \texttt{amsart} preamble that restates the article's own declarations and macros used by the retained lines, namely the environments \texttt{definition}, \texttt{lemma}, \texttt{theorem} on separate counters, together with the standard packages those lines need.  The retained environments print in this document as Definition 1, Lemma 1, Theorem 1 in order of appearance, whereas the article prints the same items with the numbering of its own full document.  The complete native source of the article as distributed in the arXiv e-print for this exact version is bundled unchanged as \texttt{sources/195253/source0.tex}.
\begin{definition}[Hyperaddition on $\Hnum$]\label{def:H-hyperadd}
Define a map
\[
  \op : \Hnum\times\Hnum
  \longrightarrow \mathcal{P}(\Hnum)\setminus\{\emptyset\}
\]
by the following rules.

\smallskip\noindent
\textbf{(1) Identity.} For all $x\in\Hnum$,
\[
  0\op x = x\op 0 := \{x\}.
\]

\smallskip\noindent
\textbf{(2) Same sign, excluding $\Lambda$.} For $\sigma\in\{+,-\}$ and
$a,b\in\Rp$,
\[
  (\sigma,a)\op(\sigma,b)
  := \{(\sigma,a+b)\}.
\]

\smallskip\noindent
\textbf{(3) Opposite real signs.} For $a,b\in\Rp$,
\[
  (+,a)\op(-,b)
  :=
  \begin{cases}
    \{(+,a-b)\}, & a>b,\\[2pt]
    \{0\},       & a=b,\\[2pt]
    \{(-,b-a)\}, & a<b.
  \end{cases}
\]
and $(-,b)\op(+,a)$ is defined by commutativity.

\smallskip\noindent
\textbf{(4) Interaction with $\Lambda$.} For $a,b\in\Rp$,
\begin{align*}
  (+,a)\op(\Lambda,b)
  &= (\Lambda,b)\op(+,a)
   := \{(\Lambda,a+b)\},\\[2pt]
  (-,a)\op(\Lambda,b)
  &= (\Lambda,b)\op(-,a)
   := \{(\Lambda,a+b)\}.
\end{align*}

\smallskip\noindent
\textbf{(5) $\Lambda$--$\Lambda$ sums.} For $a,b\in\Rp$,
\[
  (\Lambda,a)\op(\Lambda,b)
  :=
  \begin{cases}
    \{(\Lambda,2a),\,0\}, & a=b,\\[2pt]
    \{(\Lambda,a+b),\, (+,|a-b|),\, (-,|a-b|)\}, & a\ne b.
  \end{cases}
\]
\smallskip\noindent
In all cases not explicitly covered above, $\op$ is defined by
commutativity:
\[
  x\op y := y\op x.
\]
\end{definition}

\begin{lemma}[Explicit form of scalar multiplication]
\label{lem:scalar-explicit}
Let $t\in\R$ and $x\in\Hnum$.
\begin{enumerate}
  \item If $t=0$, then $t\cdot x = 0$ for all $x\in\Hnum$.
  \item If $t>0$ and $x=0$, then $t\cdot x=0$.  If
        $t>0$ and $x=(\sigma,a)$ with $\sigma\in\Snon$, $a>0$, then
  \[
    t\cdot(\sigma,a) = (\sigma,ta).
  \]
  \item If $t<0$ and $x=0$, then $t\cdot x=0$.  If
        $t<0$ and $x=(\sigma,a)$ with $\sigma\in\Snon$, $a>0$, then
  \[
    t\cdot(\sigma,a) =
    (\sigma',\,|t|a),
  \]
  where
  \[
    \sigma' =
    \begin{cases}
      -,       & \sigma=+,\\
      +,       & \sigma=-,\\
      \Lambda, & \sigma=\Lambda.
    \end{cases}
  \]
\end{enumerate}
\end{lemma}

\begin{theorem}[Distributivity for real scalars]
\label{thm:scalar-Rge0}
Let $t\in\R$ and $x,y\in\Hnum$.  Then
\[
  t\cdot(x\op y) = (t\cdot x)\op(t\cdot y)
\]
as subsets of $\Hnum$.  In particular, for each $t\in\R$ the map
$x\mapsto t\cdot x$ is an endomorphism of the additive hypermagma
$(\Hnum,\op)$.
\end{theorem}

\begin{proof}
If $t=0$, then $\iemb(t)=0$ and $0\cdot z=0$ for all $z\in\Hnum$, so
\[
  0\cdot(x\op y) = \{0\}.
\]
On the other hand, $0\cdot x=0\cdot y=0$, so
\[
  (0\cdot x)\op(0\cdot y) = 0\op 0 = \{0\}.
\]
Thus the identity holds for $t=0$.

Now assume $t>0$.  By Lemma~\ref{lem:scalar-explicit}, for every
$x=(\sigma,a)$ we have $t\cdot x=(\sigma,ta)$: scalar multiplication by
$t>0$ preserves signs and scales magnitudes by $t$.  We must show that
for each pair $x,y$ the hyper–sum of the scaled elements coincides with
the scaled hyper–sum of the original elements:
\[
  (t\cdot x)\op(t\cdot y)
  = \{t\cdot z : z\in x\op y\}.
\]

If $x=0$ or $y=0$, then $x\op y$ is a singleton by
Definition~\ref{def:H-hyperadd}, and the equality is immediate from
the fact that $t\cdot 0=0$ and $0$ is a neutral element.

If $x,y\ne0$, write $x=(\sigma,a)$, $y=(\tau,b)$ with
$\sigma,\tau\in\Snon$, $a,b>0$.  We analyse the four cases in the
definition of $\op$.

\smallskip\noindent
\emph{Case 1: Same real sign.}  Suppose $\sigma=\tau\in\{+,-\}$.  Then
by clause (2),
\[
  x\op y = \{(\sigma,a+b)\}.
\]
Hence
\[
  t\cdot(x\op y)
  = \{t\cdot(\sigma,a+b)\}
  = \{(\sigma,t(a+b))\}.
\]
On the other hand,
\[
  t\cdot x = (\sigma,ta),\quad
  t\cdot y = (\sigma,tb),
\]
so
\[
  (t\cdot x)\op(t\cdot y)
  = (\sigma,ta)\op(\sigma,tb)
  = \{(\sigma,ta+tb)\}
  = \{(\sigma,t(a+b))\}.
\]
Thus the two sets coincide.

\smallskip\noindent
\emph{Case 2: Opposite real signs.}  Suppose $\{\sigma,\tau\}=\{+,-\}$,
say $\sigma=+$, $\tau=-$; the other ordering is symmetric.  Then
\[
  x\op y
  = (+,a)\op(-,b)
  =
  \begin{cases}
    \{(+,a-b)\}, & a>b,\\[2pt]
    \{0\},       & a=b,\\[2pt]
    \{(-,b-a)\}, & a<b.
  \end{cases}
\]
If $a>b$, then $a-b>0$ and
\[
  t\cdot(x\op y)
  = \{t\cdot(+,a-b)\}
  = \{(+,t(a-b))\}.
\]
On the other hand,
\[
  t\cdot x = (+,ta),\quad
  t\cdot y = (-,tb),
\]
and since $ta>tb$ we have
\[
  (t\cdot x)\op(t\cdot y)
  = (+,ta)\op(-,tb)
  = \{(+,ta-tb)\}
  = \{(+,t(a-b))\}.
\]
The cases $a=b$ and $a<b$ are handled similarly, using the fact that
multiplication by $t>0$ preserves the inequalities $a>b$, $a=b$,
$a<b$.  In each subcase the two sets agree.

\smallskip\noindent
\emph{Case 3: Interaction with $\Lambda$.}  Suppose, for example,
$\sigma=+$, $\tau=\Lambda$; the other possibilities
$\{\sigma,\tau\}=\{-,\Lambda\}$ are analogous.  Then
\[
  x\op y
  = (+,a)\op(\Lambda,b)
  = \{(\Lambda,a+b)\}.
\]
Hence
\[
  t\cdot(x\op y)
  = \{t\cdot(\Lambda,a+b)\}
  = \{(\Lambda,t(a+b))\}.
\]
On the other hand,
\[
  t\cdot x = (+,ta),\quad
  t\cdot y = (\Lambda,tb),
\]
so
\[
  (t\cdot x)\op(t\cdot y)
  = (+,ta)\op(\Lambda,tb)
  = \{(\Lambda,ta+tb)\}
  = \{(\Lambda,t(a+b))\}.
\]
Thus the equality holds.

\smallskip\noindent
\smallskip\noindent
\emph{Case 4: $\Lambda$--$\Lambda$ sums.}  Suppose $\sigma=\tau=\Lambda$.
Then by clause (5) of Definition~\ref{def:H-hyperadd},
\[
  x\op y
  = (\Lambda,a)\op(\Lambda,b)
  =
  \begin{cases}
    \{(\Lambda,2a),\,0\}, & a=b,\\[2pt]
    \{(\Lambda,a+b),\, (+,|a-b|),\, (-,|a-b|)\}, & a\ne b.
  \end{cases}
\]
If $a=b$, then
\[
  t\cdot(x\op y)
  = \{(\Lambda,2ta),\,0\}
  = (\Lambda,ta)\op(\Lambda,ta)
  = (t\cdot x)\op(t\cdot y).
\]
If $a\ne b$, then $ta\ne tb$ and
\[
  t\cdot(x\op y)
  = \{(\Lambda,t(a+b)),\, (+,t|a-b|),\, (-,t|a-b|)\},
\]
while
\[
  (t\cdot x)\op(t\cdot y)
  = (\Lambda,ta)\op(\Lambda,tb)
  = \{(\Lambda,t(a+b)),\, (+,|ta-tb|),\, (-,|ta-tb|)\},
\]
and $|ta-tb|=t|a-b|$ for $t>0$.  Thus the equality holds.

Thus distributivity holds in all four cases.

Since every pair $(x,y)$ falls into one of these cases (or the trivial
case with a zero term), the equality
$t\cdot(x\op y)=(t\cdot x)\op(t\cdot y)$ holds for all $x,y\in\Hnum$
whenever $t>0$.  Together with the $t=0$ case, this settles all $t\ge0$.

Finally, assume $t<0$ and write $t=-s$ with $s>0$.  Let
$\iota:\Hnum\to\Hnum$ be the involution that fixes $0$ and $\Lambda$ and
swaps $+$ and $-$:
\[
  \iota(0)=0,\quad \iota(+,a)=(-,a),\quad \iota(-,a)=(+,a),\quad \iota(\Lambda,a)=(\Lambda,a).
\]
A direct inspection of Definition~\ref{def:H-hyperadd} shows that $\iota$
is an automorphism of $(\Hnum,\op)$, i.e.\ $\iota(x\op y)=\iota(x)\op\iota(y)$.
Moreover, Lemma~\ref{lem:scalar-explicit} gives $t\cdot x=\iota(s\cdot x)$ for all
$x\in\Hnum$.  Therefore
\[
  t\cdot(x\op y)=\iota\bigl(s\cdot(x\op y)\bigr)
  =\iota\bigl((s\cdot x)\op(s\cdot y)\bigr)
  =(t\cdot x)\op(t\cdot y),
\]
where we used the $s>0$ case and the fact that $\iota$ is applied elementwise
to subsets.  This completes the proof for all $t\in\R$.
\end{proof}

\end{document}
